\documentclass[conference]{IEEEtran}
\usepackage{graphicx}
\usepackage{amsmath}
\usepackage{booktabs}
\usepackage{url}
\usepackage{array}
\usepackage{tabularx}
\title{SudokuScan: From a Photograph of a Printed Puzzle to a Digital 9$\times$9 Grid}
\author{\IEEEauthorblockN{[Your Name]}\IEEEauthorblockA{Department of Computer Engineering\\MEF University\\Student no.: [Student Number]}}
\begin{document}
\maketitle

\begin{abstract}
[Write a concise summary using your actual personal-dataset results. Do not claim accuracy or robustness before measuring it.]
\end{abstract}

\section{Introduction}
[Problem, challenges, approach, dataset summary and measured headline results.]

\section{Corner Detection}
\subsection{Harris Parameters and Candidate Suppression}
[Give blockSize, ksize, k, response threshold and NMS radius. Explain their roles.]
\subsection{Automatic Outer-Corner Selection}
[Explain region proposals, how selected coordinates are guaranteed to originate from Harris candidates, candidate scoring and sub-pixel refinement.]
\subsection{Parameter Study}
[Insert a table/plot from at least six personal images. Include candidate counts and corner errors for the tested settings.]
\begin{figure}[t]
\centering
\includegraphics[width=\linewidth]{figures/harris_parameter_study.png}
\caption{Harris parameter study. Replace this placeholder with your measured plot.}
\end{figure}
\subsection{Failure Cases}
[Discuss clutter, shadows, rotation, partial borders and other observed failures.]

\section{Rectification}
\subsection{Affine and Projective Models}
[Explain A3, A4 and H4 and their degrees of freedom.]
\subsection{Residuals}
[Report A3 unused-BR error, A4 four-corner RMS residual and H4 mapping residual.]

\section{Affine Versus Homography}
[Define the obliqueness measure and show results for at least eight images, near-frontal through strongly oblique.]
\begin{table}[t]
\centering
\caption{Rectification comparison on the personal dataset. Replace every placeholder with measured values.}
\begin{tabular}{lrrrr}
\toprule
Image & Obliqueness & A3 error & A4 RMS & H4 RMS\\
\midrule
[photo 1] & [ ] & [ ] & [ ] & [ ]\\
[photo 2] & [ ] & [ ] & [ ] & [ ]\\
\bottomrule
\end{tabular}
\end{table}
[Compare grid-line deviations and OCR cell accuracy for A3/A4/H4. Relate your evidence to weak perspective.]
\begin{figure}[t]
\centering
\includegraphics[width=\linewidth]{figures/corner_residual_vs_obliqueness.png}
\caption{Affine and projective corner residuals against obliqueness.}
\end{figure}
\begin{figure}[t]
\centering
\includegraphics[width=\linewidth]{figures/grid_straightness_vs_obliqueness.png}
\caption{Grid-line deviation against obliqueness.}
\end{figure}
\begin{figure}[t]
\centering
\includegraphics[width=\linewidth]{figures/ocr_accuracy_vs_obliqueness.png}
\caption{OCR cell accuracy against obliqueness.}
\end{figure}
\begin{figure}[t]
\centering
\includegraphics[width=\linewidth]{figures/grid_spacing_vs_obliqueness.png}
\caption{Grid-line spacing deviation against obliqueness.}
\end{figure}
\begin{figure}[t]
\centering
\includegraphics[width=\linewidth]{figures/most_oblique_rectifications.png}
\caption{A3, A4 and H4 on the most oblique personal image.}
\end{figure}

\section{Recognition}
\subsection{Cell Margin and Empty-Cell Decision}
[Explain the margin and explicit ink/component threshold. Emptiness must be decided before OCR.]
\subsection{OCR Configuration and Cell Preparation}
[State engine/version, PSM, whitelist, binarization, component filtering, centring and upscaling. Compare raw and prepared cells using ground truth.]
\subsection{Recognition Results}
[Report overall cell accuracy, empty-cell F1 and most frequent digit confusions. Include representative cell crops.]

\section{Results}
[Provide corner error, cell accuracy, empty-cell F1, exact-grid match and robustness, broken down by C1-C5. Include the development set only if your course-provided annotations are available to you.]

\section{Failure Analysis}
[At least four genuine failures. For each: input crop, Harris map, debug image, failing stage, cause and a proposed correction.]

\section{Conclusion and Limitations}
[Evidence-based conclusion: when is affine sufficient, when is H4 needed, and what remains unreliable?]

\appendices
\section{AI-Use Disclosure}
The use described below must be checked and edited to match the student's actual use. Add every other AI tool used and remove inaccurate details.
\begin{table*}[t]
\centering
\caption{AI-use disclosure (Appendix C).}
\begin{tabularx}{\textwidth}{p{2.2cm} p{2.4cm} X X}
\toprule
Tool & Purpose / stage & What was asked (summary) & What was kept, changed or rejected \\
\midrule
ChatGPT (OpenAI) & Project planning, Python implementation, code review/fixes, tests, Windows installer/launcher, documentation and optional B1-B4 & Create and then audit a project implementing the supplied COMP 482 SudokuScan specification; fix identified bugs and add one-click Windows scripts & The generated implementation and subsequent fixes are included. The student must inspect, run, understand, validate, modify and take responsibility for submitted parts. Edit this description to match actual use. \\
\bottomrule
\end{tabularx}
\end{table*}

\end{document}
